\section{Weighted consequences}\label{sec:weighted}

\begin{contract}
\textbf{Assumes:} \cref{thm:main}. \quad
\textbf{Produces:} the vertex- and arc-weighted forms of the conjecture
(\cref{cor:weighted}). \textbf{Why:} these are the forms in which the conjecture is usually applied;
the passage from the unweighted statement to them is Seacrest's, and is quoted rather than reproved.
Nothing in this section is used in the proof of \cref{thm:main}.
\end{contract}

The weighted forms come in two quantifier orders, and the difference matters. One can ask that
\emph{for every} weighting there be a vertex that does well (a max-min statement), or that there be
\emph{one} weighting for which every vertex does well (a min-max statement). Both hold, and both
follow from \cref{thm:main} through the equivalence proved by Seacrest~\cite{Seacrest2015}.

\begin{corollary}[Weighted second-neighborhood forms]\label{cor:weighted}
Let $D$ be a nonempty finite oriented graph and $A(D)$ its arc set. For a vertex weighting
$\eta:V(D)\to[0,\infty)$ write $\eta(S)=\sum_{x\in S}\eta(x)$. Then:
\begin{enumerate}[label=\textup{(\roman*)},leftmargin=1.6em,itemsep=3pt]
\item for every vertex weighting $\eta:V(D)\to[0,\infty)$ there is a vertex $v$ with
  \[
   \eta(N_1^+(v))\le\eta(N_2^+(v));
  \]
\item there is a vertex weighting $\eta:V(D)\to[0,\infty)$ with $\sum_{x\in V(D)}\eta(x)=1$ such that
  \[
   \eta(N_1^+(v))\le\eta(N_2^+(v))\qquad\text{for every }v\in V(D),
  \]
  whose entries need not all be positive;
\item for every arc weighting $w:A(D)\to[0,\infty)$, extended by $w(vs)=0$ when $vs\notin A(D)$, and
  with
  \[
   \beta_v^w(s)=\max\Bigl(\{0\}\cup\{w(us)-w(vs):\ u\in V(D),\ v\to u\to s\}\Bigr),
  \]
  there is a vertex $v$ with
  \[
   \sum_{s\in V(D)}\beta_v^w(s)\ \ge\ \sum_{u\in N_1^+(v)}w(vu).
  \]
\end{enumerate}
\end{corollary}

\inwords{Weighted by nonnegative numbers, either every weighting admits a good vertex, or one
weighting makes every vertex good; and the same holds for weights on arcs, with a two-step maximum
in place of the second neighborhood.}

\begin{proof}
\emph{(iii)} is Seacrest's arc-weighted conjecture, which his Proposition~2 shows to be equivalent to
the unweighted conjecture~\cite{Seacrest2015}; \cref{thm:main} is the unweighted conjecture, so (iii)
holds.

\emph{(i) follows from (iii).} Given $\eta$, set $w(us)=\eta(s)$ for every arc $us$. Then for
$v,s\in V(D)$,
\[
 \beta_v^w(s)=\begin{cases}\eta(s),&s\in N_2^+(v),\\ 0,&s\notin N_2^+(v),\end{cases}
\]
as \cref{fig:beta} illustrates: if $v\to s$ then every eligible difference
$w(us)-w(vs)=\eta(s)-\eta(s)$ is $0$, and if $s$ is at exact distance two from $v$ then $w(vs)=0$ and
the maximum is $\eta(s)$. So the two sides of (iii) become $\eta(N_2^+(v))$ and $\eta(N_1^+(v))$,
which is (i).

\emph{(ii) follows from (i) by duality.} Let $M_D$ be the matrix with rows and columns indexed by
$V(D)$, with entry $1$ at $(v,s)$ when $s\in N_2^+(v)$, entry $-1$ when $s\in N_1^+(v)$, and entry
$0$ otherwise. Reversing all arcs transposes this matrix, $M_{\overleftarrow D}=M_D^{\mathsf T}$,
because reversing swaps the roles of first and second out-neighborhoods. By Farkas' lemma in the form
used by Seacrest~\cite[Theorem~1]{Seacrest2015}, exactly one of the following holds: there is a
nonzero $\eta\ge0$ with $M_D\eta\ge0$, or there is $\xi\ge0$ with $M_D^{\mathsf T}\xi<0$ entrywise.
The second alternative is impossible: it would say that in the reversed graph $\overleftarrow D$ every
vertex $v$ has $\xi(N_1^+(v))>\xi(N_2^+(v))$, contradicting (i) applied to $\overleftarrow D$ with the
weighting $\xi$. So the first alternative holds; normalizing $\eta$ to total weight $1$ gives (ii).
\end{proof}

\begin{figure}[t]
\centering
\begin{tikzpicture}[
  v/.style={circle,draw=black!70,fill=black!8,inner sep=0pt,minimum size=5.2mm,font=\scriptsize},
  ar/.style={-{Stealth[length=1.5mm]},draw=black!75},
  lab/.style={font=\scriptsize,text=black!75},
  wl/.style={font=\scriptsize,text=blue!45!black,align=left},
  pan/.style={font=\scriptsize\itshape,text=black!55}
]
% ---------------- left: the definition
\node[v] (v) at (0,1.3) {$v$};
\node[v] (u) at (1.6,2.4) {$u$};
\node[v] (s) at (3.2,1.3) {$s$};
\draw[ar] (v) -- (u);
\draw[ar] (u) -- node[above right=0pt,lab] {$w(us)$} (s);
\draw[ar] (v) -- node[below right=0pt,lab] {$w(vs)$} (s);
\node[lab,anchor=north] at (1.6,0.55) {$\beta_v^w(s)=\max\bigl(0,\ w(us)-w(vs)\bigr)$};
\node[pan,anchor=north] at (1.6,-0.1) {left: the definition};
% ---------------- right: the specialization
\begin{scope}[shift={(7.4,0)}]
\node[v] (v2) at (0,1.3) {$v$};
\node[v] (u2) at (1.5,2.4) {$u$};
\node[v] (s3) at (3.0,1.3) {$s$};
\node[v] (t2) at (3.0,0.2) {$t$};
\draw[ar] (v2) -- (u2);
\draw[ar] (u2) -- (s3);
\draw[ar] (v2) -- (s3);
\draw[ar] (v2) -- (t2);
\node[wl,anchor=west] at (3.45,1.6) {$s\in N_2^+(v)$:\\ $\beta=\eta(s)$};
\node[wl,anchor=west] at (3.45,0.2) {$t\in N_1^+(v)$:\\ $\beta=0$};
\node[pan,anchor=north] at (1.5,-0.1) {right: $w(us)=\eta(s)$};
\end{scope}
\end{tikzpicture}

\caption{The arc weighting $\beta_v^w$ of \cref{cor:weighted}(iii), and the substitution of (i).
\emph{Left:} for each two-step path $v\to u\to s$ one compares the last arc weight $w(us)$ with the
first arc weight $w(vs)$; $\beta_v^w(s)$ is the largest excess over all intermediate vertices $u$,
and $0$ if there is none. \emph{Right:} with the constant weighting $w(us)=\eta(s)$ the excess is
$\eta(s)$ at an exact second out-neighbor and $0$ everywhere else, so the inequality of (iii) becomes
$\eta(N_2^+(v))\ge\eta(N_1^+(v))$. \emph{Why it matters:} the figure shows why the arc-weighted form
is the right common generalization: its two sides are a sum over out-arcs and a sum over two-step
excesses, and the vertex-weighted statements are the two specializations.}
\label{fig:beta}
\end{figure}

\begin{remark}[What is new here and what is quoted]\label{rem:quoted}
Only \cref{thm:main} is new. The passage from it to \cref{cor:weighted} is Seacrest's: (iii)
$\Leftrightarrow$ unweighted is his Proposition~2, and the Farkas duality between (i) and (ii) is his
Theorem~1. We quote these rather than reprove them, and we have verified that the statements quoted
are the ones used: in particular, (iii) is the form with a \emph{maximum over intermediate vertices},
not a sum over paths, and the arc weights are allowed to vanish.
\end{remark}
